\guidechapter{See exactly which digit products get added}{Convolution without new machinery.}

Every term in long multiplication has a destination place. In $23\times14$,
the $3\times4$ product contributes to the ones place. Both $3\times1$
and $2\times4$ contribute to the tens place, because one factor in each
product is a tens digit. Finally $2\times1$ contributes to the hundreds
place. The diagonal colors in the grid identify those destinations.

\explfigure{13-convolution-groups}{Convolution is addition along equal-place diagonals.}{
The two teal cells contribute $3$ and $8$ to the same coefficient, giving
$11$. The coefficient list is $[12,11,2]$, before carrying. A cell is a
specific digit product, not a generic stage of an algorithm.}

Number digit positions starting at zero. A digit in position $i$ multiplies
a digit in position $j$ into position $i+j$, because $10^i10^j=10^{i+j}$.
That is the whole reason for the rule ``add products with the same position
sum.'' The word \emph{convolution} names this rule. It is the same rule in
binary or any other radix.

\paragraph{Why change to a polynomial?}
Replace the radix ten by a variable $x$. The two inputs become
$A(x)=3+2x$ and $B(x)=4+x$. Their product is a quadratic: it can be
written $C(x)=c_0+c_1x+c_2x^2$ for three unknown coefficients.
At $x=10$ it must give the integer product. But ten is not the only useful
place to evaluate it. Easier sample points can reveal its coefficients.

\explfigure{14-evaluate-recover}{Compute values first, then recover coefficients.}{
At $x=-1,0,1$, the input values multiply to $3,12,25$. Exactly one polynomial
of degree at most two passes through those three distinct sample points.
The curves are these specific polynomials, not a metaphor for a transform.}

At zero, $C(0)=c_0=3\times4=12$. At one, $C(1)=5\times5=25$.
At minus one, $C(-1)=1\times3=3$. Thus
\[
 c_0=12,\qquad
 c_1=\frac{25-3}{2}=11,\qquad
 c_2=\frac{25+3}{2}-12=2.
\]
Subtracting the last two sampled equations isolates the odd-power term;
adding them isolates the even-power terms. Evaluating at ten now gives
$12+11\cdot10+2\cdot100=322$.

\paragraph{Where this example stops.}
This tiny interpolation example is not the paper's fast algorithm. It shows
why ``evaluate, multiply matching values, reconstruct'' can work at all.
For huge inputs, evaluation and reconstruction must themselves be fast and
well controlled. Fourier transforms choose highly structured sample points
so that those steps can share work. The next foundation chapter explains
what their patterns look like.

\takeaway{Convolution groups equal-place products. A transform provides
another way to compute the same coefficients; it does not change the product.}
